\section{Setup: one loss, two knobs, one budget}
\label{sec:setup}

\paragraph{One loss for every method.}
All contrastive methods in this paper are the same loss with a different definition of ``positive pair''.
A batch holds $B$ images; each is augmented twice, giving $2B$ views with $\ell_2$-normalised
projections $z_1,\dots,z_{2B}$. Every image $n$ carries a \emph{group id} $g_n$, and two different views are
positives exactly when their images share a group id:
\begin{equation}
\label{eq:loss}
\mathcal{L} \;=\; \frac{1}{2B}\sum_{i=1}^{2B} \frac{-1}{|P(i)|}\sum_{p\in P(i)}
\log\frac{\exp(z_i^\top z_p/\tau_i)}{\sum_{a\neq i}\exp(z_i^\top z_a/\tau_i)},
\qquad P(i)=\{p\neq i : g_{n(p)} = g_{n(i)}\}.
\end{equation}
\begin{itemize}
\item \textbf{\simclr} \citep{chen2020simclr}: no labels are used, every image is its own group ($g_n = n$), so
$P(i)$ is just the other view of the same image and \cref{eq:loss} is the NT-Xent loss.
\item \textbf{\supcon} \citep{khosla2020supcon}: every image is labeled and its group is its class
($g_n = y_n$); \cref{eq:loss} is the $\mathcal{L}^{sup}_{out}$ loss of \citet{khosla2020supcon}.
\item \textbf{\union}: labeled images use their class ($g_n=y_n$) and unlabeled images are their own group
($g_n = n$). An unlabeled image is still contrasted against every other view in the batch, so it keeps
providing \emph{negatives} and its own \emph{instance} positive; it simply never claims a class.
\end{itemize}
\union is not a new idea: it is the natural semi-supervised reading of \cref{eq:loss}, and close to
SuNCEt \citep{assran2020suncet} and to the joint SimCLR+SupCon objectives of \citet{islam2021broad}.
We use it as a \emph{probe} rather than as a proposed method: it differs from label-only \supcon only in
which images are in the pool, and from \simclr only in whether the labels are used. Two details of \union were chosen on a held-out
validation split of the \emph{training} set before any test-set number was computed (\cref{app:pilot}):
labeled anchors use $\tau_i=0.1$ and unlabeled anchors $\tau_i=0.5$ (the standard temperatures of
\supcon and \simclr on CIFAR), and 25\% of every batch is drawn from the labeled images. A variant with uniform batches was added after we saw the long-budget results and is reported only as an ablation.
We verified our implementation of \cref{eq:loss} numerically against the official \supcon code and our own
NT-Xent implementation (maximum absolute difference below $10^{-6}$, \cref{app:impl}).

\paragraph{Two knobs: labels and images.}
With a labeled fraction $f$ (10\% unless stated), the two quantities that the label-scarce setting couples
are the \emph{labels} the loss may use and the \emph{images} in the pretraining pool. We vary them
independently (\cref{tab:main}):
\begin{center}
\small
\begin{tabular}{lcc}
\toprule
 & pool = the $f$ labeled images & pool = all images \\
\midrule
labels not used & \simclr (subset) & \simclr \\
labels used     & \supcon (label-only) & \union \\
\bottomrule
\end{tabular}
\end{center}
Comparing the rows isolates the effect of the labels at a fixed pool; comparing the columns isolates the
effect of the unlabeled images at a fixed use of labels. The labeled subsets are stratified (the same
number of images per class) and nested across fractions for a given seed, so the 1\% subset is part of
the 5\% subset, and so on.

\paragraph{One budget: steps, not epochs.}
In every run we fix the number of gradient steps to $E\cdot\lfloor N/B\rfloor$, where $N=50{,}000$ is the size
of the \emph{full} training set, $B=512$ and $E=100$ unless stated. ``100 epochs'' therefore means 9,700
steps of 1,024 augmented views for every method, whatever the size of its pool; a method with a
5,000-image pool simply cycles through it ten times as often. We also report budgets from 970 to 19,400 steps
(``10'' to ``200 epochs''); \cref{sec:trap} shows why this matters.

\paragraph{Training details.}
All encoders are ResNet-18 \citep{he2016resnet} with the CIFAR stem (a $3{\times}3$ stride-1 first
convolution and no max-pooling) and a two-layer projection head with batch normalisation in the hidden
layer only, as in the official \supcon code (adding batch normalisation on the output made \supcon
collapse in our earlier experiments, \cref{app:report}). We train with SGD (momentum 0.9, weight decay
$10^{-4}$), learning rate $0.25$ for batch 512 (the official recipe's $0.5$ at batch 1024), 5\% linear
warm-up and a cosine schedule, in mixed precision with the loss computed in float32. The augmentation is
the CIFAR recipe of \citet{chen2020simclr}: random resized crop (scale $0.2$--$1$), horizontal flip, colour
jitter with strength $s=0.5$ applied with probability $0.8$, and random grayscale ($p=0.2$). We wrote
batched GPU versions of these transforms and checked each of them against \texttt{torchvision} (maximum
absolute pixel difference below $3\times10^{-7}$; \cref{app:impl}); this makes a 100-epoch ResNet-18 run take
about 13 minutes on one NVIDIA A40. Full hyper-parameters, including those of the baselines, are listed in
\cref{app:hparams}.

\paragraph{Evaluation.}
The encoder is frozen and its 512-d average-pooled features are used in four ways:
(i) \textbf{linear probe} with all 50,000 training labels, the standard protocol of
\citet{chen2020simclr}: an $\ell_2$-regularised multinomial logistic regression on standardised features,
fit with L-BFGS, with the regularisation strength picked on a stratified 20\% split of the
\emph{training} labels; (ii) \textbf{weighted $k$NN} ($k=200$, temperature $0.1$; \citealp{wu2018instance});
(iii) the same linear probe but trained \textbf{only on the labeled subset}; and (iv) \textbf{fine-tuning} of the
whole network on the labeled subset (linear-probe initialisation followed by 60 epochs of SGD over the labeled
images; \citealp{kumar2022lpft}). Protocols (i) and (ii) measure representation quality as the
self-supervised literature does; (iii) and (iv) are what a practitioner who has only $f$ of the labels can
actually do, and for the label-using methods they reuse the labels seen in pretraining. The test set is used
only for these final numbers.

\paragraph{Statistics.}
Main configurations are run with \nSeeds seeds; the seed changes the labeled subset, the initialisation
and the augmentation stream. We report mean $\pm$ standard deviation over seeds and, for every key
comparison, a 95\% confidence interval and $p$-value from a bootstrap that resamples test images
(paired: both methods are scored on the same resampled images) \emph{and} seeds, so that the interval
reflects both sources of noise \citep{efron1993bootstrap}. Exploratory sweeps use one seed and are
labeled as such.
